\section{Related Work}% \gm{We have the space to add narrative commentary in here}% 
% train model per embodiment:
% \cite{co24, ACT23, to24, xo24}
% Train and fine-tune per embodiment:
% \cite{pi_0_24, chen_conrft_2025, wzltsf25-1, kim_openvla_2024, yo24, shukor_smolvla_2025, wzltsf25, nvidia_gr00t_2025, wiut24, yo25}

% Policies Engineered for Cross-Embodiment: \cite{bogert_built_2024, yzfl24, freiberg_diffusion_2024, chen_graphmimic_nodate, lo24-1, yo25}

% Latent Space Policies:
% \cite{wo25-1, bauer_latent_2025}

% Others:
% \cite{seo_masked_2023, yang_pushing_2024, ldb25}

% Cross-embodiment learning seeks to develop control policies that function across heterogeneous robot bodies—i.e., differing morphologies, kinematics, sensors, and actuators—so that capability does not have to be relearned from scratch for each platform. train model per embodiment specialized, robot/task-specific policies that optimize in-domain performance without explicit transfer guarantees \cite{co24, ACT23, to24, xo24}. \emph{Train and fine-tune per embodiment}—generalist VLAs/flow models pretrained on diverse multi-robot data and adapted to each target via PEFT or brief online updates, balancing reuse with hardware specificity \cite{pi_0_24, chen_conrft_2025, wzltsf25-1, kim_openvla_2024, yo24, shukor_smolvla_2025, wzltsf25, nvidia_gr00t_2025, wiut24, yo25}. \emph{Policies Engineered for Cross-Embodiment}—representation alignment, shared interfaces, and cross-domain co-training that aim for one policy spanning morphologies \cite{bogert_built_2024, yzfl24, freiberg_diffusion_2024, chen_graphmimic_nodate, lo24-1, yo25}. \emph{Latent Space Policies}—robot-agnostic action/interaction spaces that unify end-effectors and enable transfer with thin adapters \cite{wo25-1, bauer_latent_2025}. \emph{Others}—complementary lines on masking, scaling across embodiments, and broad evaluation \cite{seo_masked_2023, yang_pushing_2024, ldb25}.
% Cross-embodiment in robotics addresses the fact that skills learned on one robot rarely transfer to another because morphologies, sensors, and action spaces differ; the aim is to reduce per-robot data/compute and enable policy reuse across hardware. At one end, some works train a model per embodiment—specialized robot- and task-specific policies that maximize in-domain reliability, with occasional use of intermediate interfaces like object flow to bridge human/sim data to a target platform \cite{co24, ACT23, to24, xo24}. Another line pretrains generalist VLA or flow models on heterogeneous multi-robot corpora and then fine-tunes them for each target with parameter-efficient updates or brief online RL, yielding dependable last-mile adaptation while retaining broad priors \cite{pi_0_24, chen_conrft_2025, wzltsf25-1, kim_openvla_2024, yo24, shukor_smolvla_2025, wzltsf25, nvidia_gr00t_2025, wiut24, yo25}. A complementary direction deliberately engineers morphology-robust interfaces and representations—tactile shared interfaces, gripper-agnostic encoders with geometry-aware decoders, graph abstractions, adversarial decomposition, and unsupervised cross-embodiment pretraining—so a single policy or pretrained backbone can span mismatched bodies with only thin adapters \cite{bogert_built_2024, yzfl24, freiberg_diffusion_2024, chen_graphmimic_nodate, lo24-1, yo25}. Finally, latent-space policies learn a robot-agnostic interaction/action latent and pair it with lightweight embodiment decoders or hand descriptors, enabling one controller to serve multiple end-effectors with minimal retraining \cite{wo25-1, bauer_latent_2025}.
% Cross embodiment addresses the gap where a skill trained on one robot embodiment often fails on another. This is because embodiment, sensing, and action spaces 
% differ. A well established line of research is to train a model per embodiment, prioritizing reliable in domain performance and, when leveraging external data,  and occasionally using intermediate interfaces  to bridge datasets and  \cite{co24, ACT23, to24, xo24}. Other works pretrain generalist VLA or flow matching models on heterogeneous multi robot dataset \cite{pi_0_24, chen_conrft_2025, wzltsf25-1, kim_openvla_2024, yo24, shukor_smolvla_2025, wzltsf25, nvidia_gr00t_2025, wiut24, yo25}. They then fine tune each model for the target robot with parameter efficient updates or brief online RL, reusing broad priors while delivering dependable last mile performance. Another direction engineers embodiment robust interfaces and representations. Examples include tactile shared interfaces, gripper agnostic encoders with geometry aware decoders, graph based video abstractions, adversarial decomposition, and unsupervised cross embodiment pretraining, enabling one policy or a pretrained backbone to span mismatched robots with only thin adapters \cite{bogert_built_2024, yzfl24, freiberg_diffusion_2024, chen_graphmimic_nodate, lo24-1, yo25}. Finally, latent space policies learn a robot-agnostic interaction or action latent and pair it with lightweight embodiment decoders or hand descriptors, enabling one controller to serve multiple end effectors with minimal retraining \cite{wo25-1, bauer_latent_2025}.





% \gm{add papers I want to cite}
% \subsection{Embodiment Failure Recovery}

% Robots operating over long horizons in uncertain environments must cope with degradation to their physical embodiment, including joint locking, reduced range, and dynamic instability. Analytical approaches have addressed this by explicitly modifying motion planning and control algorithms to maintain function under such impairments. Techniques such as self-motion manifold intersection planning~\cite{xie2021maximizing}, fail-safe reachability analysis~\cite{porges2021failsafe}, and redundancy reconfiguration via closed-loop inverse kinematics~\cite{rayankula2021fault} aim to guarantee task success despite failures. Others use optimization-based control to exclude failed joints~\cite{yang2024qp}, derive inverse kinematics for specific locked-joint cases~\cite{wang2021ik}, or adapt control laws for stability under chaotic degradation~\cite{khan2023chaos}. These methods provide robust, explainable guarantees, but are often limited to single failure modes, require significant manual modeling effort, and do not scale easily to unseen configurations. \gm{fix this sentence, maybe make it two sentences, maybe incoproate/sprinkle the limitations throughout?}

% Design-time strategies aim to increase fault tolerance by optimizing kinematic structure itself~\cite{almarkhi2020design}, while task-specific gait reconfiguration for legged systems~\cite{xi2025gait} allows continued locomotion with partial joint paralysis \gm{note how a lot of these are on walking and some on point to point, but none on task solutions}. Collectively, these works show that it is possible to analytically adapt to embodiment failures when the robot model is known and the failure mode is anticipated. Our proposed approach takes a contrasting view: instead of hand-crafting failure-specific solutions, we use a data-driven generative model trained on diverse failure data to synthesize feasible, compensatory trajectories at inference time, potentially generalizing across multiple failure types and robots.

% \subsection{Learning-Based Failure Recovery}

% A growing body of learning-based work aims to overcome the limitations of handcrafted methods by training robust, generalizable policies. Reinforcement learning has been used to adapt to joint-level failures via adversarial training~\cite{yang2021adversarial}, partial observability~\cite{pham2024adaptive}, and random joint masking during locomotion~\cite{kim2024quadrupedal}. These strategies have proven capable of recovering from severe physical impairments, particularly in quadrupeds and manipulators. Policy architectures such as FT-Net~\cite{luo2023ftnet} and fault-conditioned networks enable recovery behaviors to activate automatically during failure. \gm{todo: weaknesses}

% Other strategies leverage curriculum learning to gradually expose the robot to increasing impairment severity~\cite{okamoto2022acdr}, or build behavior repertoires via quality-diversity search that can rapidly switch to a viable alternative after damage~\cite{cully2015qd, allard2023hierarchical} \gm{double check the exact method cully uses}. While these approaches often demonstrate strong performance, they typically rely on task-specific training pipelines and either explicit retraining or policy selection at runtime. In contrast, our diffusion-based model learns a generative distribution of full trajectories over failure modes, allowing flexible inference-time adaptation without online optimization or policy switching.
\paragraph{Failure-Aware Robot Control}

Robots operating over long durations in uncertain environments must be able to continue functioning even when experiencing physical degradation. Joint-level failures, such as locking, reduced range of motion, or diminished velocity, introduce complex, nonlinear constraints that disrupt standard motion planning approach. Classical methods attempt to address these challenges through explicit algorithmic adaptation: self-motion manifold planning~\cite{xie2021maximizing}, fail-safe reachability analysis~\cite{porges2021failsafe}, and redundancy exploitation via inverse kinematics~\cite{rayankula2021fault}. In assuming a subset of possible failures, these approaches offer guarantees on post-failure ability but they are unable to extend to arbitrary failure conditions. Other works adapt control laws to handle specific joint failures~\cite{yang2024qp, wang2021ik, khan2023chaos}, but are often limited in generality, scale poorly with the combinatorial space of failure types, and rely heavily on manual modeling. While effective in narrow domains, these handcrafted solutions cannot address the full complexity of real-world fail-active behavior, where failures must be handled, and task strategies adapted, on the fly.

% Learning-based strategies aim to overcome the limitations of traditional methods by training control policies that adapt to embodiment changes through experience. In reinforcement learning, damage-aware behaviors have been demonstrated via adversarial training~\cite{yang2021adversarial}, partial observability~\cite{pham2024adaptive}, and stochastic joint masking~\cite{kim2024quadrupedal}. Other approaches use curriculum learning~\cite{okamoto2022acdr} or quality-diversity search~\cite{cully2015qd, allard2023hierarchical} to build behavioral repertoires capable of failure recovery. 

Learning-based strategies address the limitations of traditional methods by training control policies that adapt to embodiment changes through experience. In reinforcement learning, damage-aware behaviors have been achieved using adversarial training \cite{yang2021adversarial}, partial observability \cite{pham2024adaptive}, and stochastic joint masking \cite{kim2024quadrupedal}. Other methods apply curriculum learning \cite{okamoto2022acdr} or quality diversity search \cite{cully2015qd, allard2023hierarchical} to develop behavioral repertoires for failure recovery.
These techniques have shown success in recovering locomotion or point-to-point reaching under known embodiments, but often require task-specific training loops, explicit policy switching, or runtime optimization. 
% More importantly, no existing method generalizes across multiple manipulation primitives or failure configurations without retraining. 
In contrast, our approach trains a generative diffusion model over full trajectories, enabling flexible, zero-shot adaptation to joint degradation and task shifts without policy switching or fine-tuning.

\paragraph{Diffusion Models for Trajectory Generation}
Diffusion models offer key advantages for trajectory generation in fail-active robotic systems, especially when embodiment changes unpredictably. Compared to imitation or reinforcement learning, diffusion models (1) capture complex, multimodal action distributions essential for recovery behaviors \cite{chi2023diffusion}, (2) provide stable training, outperforming alternatives like energy-based models \cite{zhong2023edmp}, and (3) enable conditional generation on structured inputs such as goals or embodiment states \cite{yang2023dimsam}. These features allow robots to adapt to joint failures without explicit retraining.

Empirically, diffusion models excel in robotic tasks like navigation, manipulation, and object rearrangement \cite{sridhar_nomad_2023}, outperforming reinforcement learning and behavior cloning in long-horizon, contact-rich tasks \cite{dasari_ingredients_2024}. Advances in diffusion transformers, like adaptive normalization and efficient tokenization, further enhance their suitability for continuous control \cite{dasari_ingredients_2024}. Crucially, diffusion supports online adaptation via trajectory reconditioning, enabling flexible fail-active responses without policy switching \cite{zhong2023edmp}, making them ideal for handling joint failures that reshape the motion space.


% Diffusion models offer key advantages for trajectory generation in fail-active robotic systems, particularly where the robot's embodiment can change unpredictably. Compared to conventional imitation or reinforcement learning methods, diffusion models (1) capture complex, multimodal action distributions essential for diverse recovery behaviors \cite{chi2023diffusion}, (2) maintain stable and reliable training dynamics,even relative to alternative generative methods like energy-based models \cite{zhong2023edmp}, and (3) enable conditional generation on structured inputs such as goals or embodiment states \cite{yang2023dimsam}. These properties are critical when robots must flexibly adapt to new constraints imposed by joint failures without explicit retraining or reprogramming.

% Beyond their architectural strengths, diffusion models have demonstrated strong empirical performance in real-world robotic domains including navigation, object rearrangement, and manipulation \cite{sridhar_nomad_2023}. Notably, they outperform reinforcement learning and behavior cloning in long-horizon, contact-rich tasks, where precise, multimodal reasoning is required \cite{dasari_ingredients_2024}. Recent advances in diffusion transformer architectures, such as adaptive normalization and efficient tokenization, further improve their viability for continuous control \cite{dasari_ingredients_2024}. Critically for failure recovery, diffusion models naturally support sampling-based online adaptation: they can recondition trajectories dynamically as embodiment constraints evolve, enabling flexible fail-active behavior without policy switching \cite{zhong2023edmp}. This makes diffusion particularly well-suited for real-world scenarios where joint-level failures redefine the feasible motion space and require rapid synthesis of new, task-complete trajectories.
